\section{Tiên đề là ràng buộc, không phải xuất phát điểm}
\label{sec:ch06-tien-de}

\index{tiên đề}%
\index{tính tuyến tính}%
Câu hỏi trước phương pháp cho thấy một tiên đề chỉ có ý nghĩa sau khi ta biết
nó đang cấm điều gì. Bài báo xét ba điều kiện quen thuộc: tính đầy đủ với một
baseline cố định, độ nhạy và \tn{tính tuyến tính}{linearity}~\autocite{p17firstprinciples}.
Chúng hấp dẫn vì mỗi điều kiện loại một lỗi rõ ràng: bỏ sót chênh lệch đầu ra,
gán điểm cho đặc trưng không ảnh hưởng, hoặc xử lý tổng hai mô hình khác với
tổng hai lời giải thích và nhân mô hình với một hằng số mà không nhân lời giải
thích với hằng số ấy. Chú ý rằng độ nhạy ở đây theo nghĩa bài báo dùng: một
đặc trưng không bao giờ ảnh hưởng tới đầu ra thì nhận attribution bằng không ở
mọi đầu vào. Đó là điều kiện thứ hai mà bài Integrated Gradients cũng nêu bên
cạnh điều kiện chương~\ref{ch:attribution-theo-gradient} đã gọi bằng tên ấy,
và hai điều kiện không trùng nhau.

\begin{theorem}
\label{thm:ch06-ba-tien-de}
Với $f$ khả vi liên tục hai lần trên $[0,1]^d$, mọi attribution thỏa cả ba
điều kiện đều bằng gradient của $f$ tại một điểm, nhân từng thành phần với
hiệu giữa đầu vào và baseline, tức dạng Gradient$\mathbin{\times}$Input tính
theo baseline, cộng với attribution của phần dư Taylor bậc hai của $f$ quanh
điểm ấy. Nếu attribution còn Lipschitz liên tục theo mô hình, nghĩa là đổi mô
hình một chút thì attribution chỉ đổi một chút, phần dư ấy bị chặn bởi tích của
hằng số Lipschitz, số chiều $d$ và độ cong lớn nhất của
$f$~\autocite{p17firstprinciples}.
\end{theorem}

Khi mô hình tuyến tính, phần dư trong định lý~\ref{thm:ch06-ba-tien-de} biến
mất. Với mạng ReLU, bài báo nêu con đường xấp xỉ qua các hàm trơn để nối
kết quả này với lớp mô hình từng đoạn affine~\autocite{p17firstprinciples}.

Kết quả không nói Gradient$\mathbin{\times}$Input là một phương pháp tồi. Nó
nói bộ tiên đề này để lại ít tự do hơn trực giác ban đầu gợi ra: ai muốn một
khái niệm attribution rộng hơn Gradient$\mathbin{\times}$Input thì không thể
giữ cả ba điều kiện cùng lúc. Nới một tiên đề hay thêm một tiên đề khác cũng
không thoát được ràng buộc ấy. Bài báo nêu hai ví dụ: đòi tính đầy đủ với mọi
baseline thay vì một baseline cố định thì dẫn tới mâu thuẫn, và các cách phát
biểu độ nhạy chỉ trên một lân cận của đầu vào thay vì trên toàn không gian mà
bài báo thử cũng vậy~\autocite{p17firstprinciples}.
Do đó, chương này không cố tìm thêm một danh sách tiên đề. Nó giữ lại đúng một
điều kiện là tính tuyến tính, và đổi hướng sang một đối tượng mà attribution
có thể được định nghĩa trực tiếp.
